% % % % % % % % % %%% Dichiarazione dei pacchetti standard.
\usepackage[italian]{babel}
\usepackage[T1]{fontenc} 
\usepackage[utf8]{inputenc}
\usepackage{amsfonts,amsmath,amssymb}
\usepackage{fancyvrb}
\usepackage{xcolor}
\usepackage[none]{hyphenat}
\usepackage{fancyhdr}
\usepackage{graphicx}
\usepackage{float}
\usepackage{smartdiagram}
\usepackage{subfigure}
\usepackage{tikz}


\definecolor{mygreen}{rgb}{.125,.5,.25}
\definecolor{bleudefrance}{rgb}{0.19, 0.55, 0.91}
%%% Personalizzazione del layout---articolata su cinque livelli.
%\usetheme{Madrid}        % layout complessivo.
%\usetheme{boxes}
%\usetheme{Pittsburgh}
%\usetheme{Median}
%\usetheme{Rochester}
%\usetheme{Pittsburgh}
%\usetheme{Madrid}        % layout complessivo.
%\definecolor{mygreen}{rgb}{.125,.5,.25}
%\definecolor{bleudefrance}{rgb}{0.19, 0.55, 0.91}
%\definecolor{ceruleanblue}{rgb}{0.16, 0.32, 0.75}
%\definecolor{bondiblue}{rgb}{0.0, 0.58, 0.71}
\usetheme{Median}
% \usetheme{Casalbo}
%\usecolortheme[named=ceruleanblue]{structure}
%\useinnertheme{circles} % layout interno. (circles rectangles rounded inmargin)
%\useoutertheme{smoothbars}  % layout esterno. (infolines miniframes smoothbars sidebar
%\usecolortheme{crane} % schema di colori.
%\usefonttheme{structurebold}  %schema dei font.
%\setbeamercovered{dynamic}
%\usecolortheme[named=blue]{structure}
%\usecolortheme{seahorse}
%\useinnertheme{circles} % layout interno. (circles rectangles rounded inmargin)
%\useoutertheme{sidebar}  % layout esterno. (infolines miniframes smoothbars sidebar
%\usecolortheme{crane} % schema di colori.
%\usefonttheme{structurebold}  % schema dei font.
%\usefonttheme{professionalfonts}
%\setbeamercovered{dynamic}
% Inutile dire che se volete tutti i default, potete risparmiarvi gli ultimi
% quattro comandi. 
\theoremstyle{definition}
\newtheorem{definizione}{Definizione}


\newcommand*\tick{\item[\checkmark]}
\newcommand*\fail{\item[\XSolidBrush]}
\renewcommand{\vec}{\boldsymbol}